\documentclass[10pt,a4paper]{amsart}
\usepackage[margin=19mm]{geometry}
\usepackage{amsmath,amssymb,mathtools}
\usepackage[scr=rsfs,cal=euler]{mathalfa}
\usepackage{xfrac}
\usepackage[unicode,hidelinks,bookmarks=false]{hyperref}
\newcommand{\ie}{i.e.\ }
\newcommand{\cf}{cf.\ }
\newcommand{\resp}{respectively\ }
\newcommand{\Subsection}{\subsection}
\providecommand{\nobreakdash}{\nobreak\hbox{-}}
\setlength{\emergencystretch}{2em}
\multlinegap0pt
\usepackage{bm}
\usepackage{bbm}
\usepackage{dsfont}
\usepackage{xspace}
\usepackage{cancel}
\usepackage{mathtools}
\usepackage{amsthm}
\makeatletter
\@ifundefined{twr}{\newtheorem{twr}{Twr}}{}
\@ifundefined{lem}{\newtheorem{lem}[twr]{Lemma}}{}
\makeatother
\DeclareMathOperator{\conv}{conv}
\DeclareMathOperator{\ext}{ext}
\DeclareMathOperator{\im}{Im}
\DeclareMathOperator{\lin}{lin}
\DeclareMathOperator{\re}{Re}
\title[$2$-strong uniqueness of a best approximation and of minimal]{$2$-strong uniqueness of a best approximation and of minimal projections in complex polytope norms and their duals}
\author{Tomasz Kobos, Grzegorz Lewicki}
\date{Source: arXiv:2504.03542v2 (CC BY 4.0)}
\begin{document}
\maketitle

\noindent\emph{Source and licence.} $2$-strong uniqueness of a best approximation and of minimal projections in complex polytope norms and their duals. Version arXiv:2504.03542v2 (CC BY 4.0), distributed under the Creative Commons Attribution 4.0 International licence, as recorded in the arXiv licence metadata for this exact version and shown on the abstract page https://arxiv.org/abs/2504.03542v2. Full rights notice: licenses/arxiv-2504.03542-CC-BY-4.0.txt.\par\smallskip

\noindent\textit{Editorial scope.} Counted unit: the Twr whose source label is \texttt{twrpolytope} in the article (source0.tex lines 604--607, source label \texttt{twrpolytope}), delivered together with the article's own complete original proof of it (source0.tex lines 608--669, one \texttt{proof} environment of 62 lines). No mathematics is written, repaired or re-derived: statement and proof are the article's own text line for line. Required context retained uncounted: lemcompact (lines 391--396); twr2strgeneral (lines 541--556); lemreal (lines 584--587). Every definition, display and label of those lines is retained unchanged. Omitted from this package and not redistributed: 1--390, 397--540, 557--583, 588--603, 670--943. Nothing inside a retained excerpt is omitted or altered: every retained body is delivered exactly as the article's lines are written, so no omission receipt of any kind accompanies this package. Citations: the retained excerpts carry no citation command at all, so no bibliography entry of the article is transcribed and no reference is redistributed. Reference adaptation: none was needed and none was made; every reference inside the delivered unit resolves inside it, so no unresolved reference or citation is produced. Printed slips kept literal: no printed word, symbol, hypothesis or number of the counted statement or of its proof was altered. Layout and class: the article's own class is \texttt{amsart} with packages including amsmath, amsthm, amscd, amsfonts, amssymb, graphicx, color, hyperref, enumerate; the wrapper is the lane's pinned \texttt{amsart} editorial wrapper at 10pt on A4, which restates the theorem-like environments the retained text needs and adds the article's own macro definitions that the retained text needs (\texttt{\textbackslash conv}, \texttt{\textbackslash ext}, \texttt{\textbackslash im}, \texttt{\textbackslash lin}, \texttt{\textbackslash re}), with extra wrapper packages (bm, bbm, dsfont, xspace, cancel, mathtools). The delivered numbering is the unit's own. Assets: no retained span contains a figure environment, a graphics-inclusion command, a PDF-page inclusion, a table or a TikZ picture; no image file is copied into this package, the package declares no asset dependency and no image file is redistributed with it. Licence: version arXiv:2504.03542v2 of this article is distributed under the Creative Commons Attribution 4.0 International licence, as recorded in the arXiv licence metadata for this exact version and shown on the abstract page https://arxiv.org/abs/2504.03542v2. A full rights notice is delivered at licenses/arxiv-2504.03542-CC-BY-4.0.txt. Characters: every retained excerpt is pure ASCII, so the delivered engine, XeLaTeX with its default Latin Modern text font, reports no missing character.
\begin{lem}
\label{lemcompact}
Let $X$ be a normed space over $\mathbb{K}$ and let $Y \subseteq X$ be its finite-dimensional linear subspace. If $x \in X \setminus Y$ is a vector such that $0$ is a unique best approximation for $x$, then for every $d>0$ there exists a constant $r_d>0$ such that
$$\|x-y\| \geq \|x\| + r_d\|y\|$$
for every $y \in Y$ with $\|y\| \geq d$.
\end{lem}

\begin{twr}
\label{twr2strgeneral}
Let $X$ be a normed space over $\mathbb{C}$ and let $Y \subseteq X$ be its finite-dimensional subspace. Suppose that $x \in X \setminus Y$ is such a vector that $0$ is its unique best approximation in $Y$ and a set
$$
E(x) = \{ f \in \ext(S_{X^*}): f(x) = \|x\| \}
$$
is a finite set consisting of $N$ functionals $f_1, \ldots, f_N \in S_{X^*}$. Let $0 \leq M \leq \|x\|$ be defined as
$$
M= sup\{|f(x)|: f \in \ext(S_{X^*} )\setminus |E|(x) \},
$$
where $ |E|(x) = \{ f \in \ext(S_{X^*}): |f(x)| = \|x\| \}$ (if $\ext(S_{X^*})\setminus |E|(x)$ is an empty set, then we take $M=0$). Let $Z \subseteq X$ be a linear subspace given as $Z=  \lin(Y \cup \{x\}).$ Assume that a linear subspace $F(Z) \subseteq \mathbb{C}^N$ defined as
$$ 
F(Z) = \{ (f_1(z), \ldots ,f_N(z)): \: z \in Z\}
$$
is a real subspace of $\mathbb{C}^N$. If an inequality $M < \|x\|$ holds, then $0$ is a $2$-strongly unique best approximation for $x$ in $Y$.
\end{twr}

\begin{lem}
\label{lemreal}
Let $X=(\mathbb{C}^n, \| \cdot \|)$ be a complex normed space with a norm, which is either a complex polytope norm with a real essential system of vertices or an adjoint complex polytope norm which with a real essential system of facets. Then for any two vectors $x, y \in \mathbb{R}^n$ we have $\|x+iy\| \geq \|x\|$. 
\end{lem}

\begin{twr}
\label{twrpolytope}
Let $X=(\mathbb{C}^n, \| \cdot \|)$ be a complex normed space with a norm that is either a complex polytope norm with a real essential system of vertices or an adjoint complex polytope norm with a real essential system of facets. Suppose that $Y \subseteq X$ is a real subspace and $x \in X \setminus Y$ is a real vector, for which $y_0 \in Y$ is a unique best approximation in $Y$. Then $y_0$ is a real vector and it is a $2$-strongly unique best approximation for $x$. Additionally, if the norm $\| \cdot \|$ is an adjoint complex polytope norm, then $y_0$ is not $\alpha$-strongly unique best approximation for any $\alpha<2$.
\end{twr}

\begin{proof}
Since $Y$ is a real subspace we have $\re(y_0), \im(y_0) \in Y$ and Lemma \ref{lemreal} implies that
$$\|x-y_0\|=\|(x-\re(y_0)) + i\im(y_0)\| \geq \|x-\re(y_0)\|,$$
which shows that the distance of $\re(y_0) \in Y$ to $x$ is not greater than of $y_0$. Since $y_0$ is a unique best approximation for $x$ it follows that $y_0=\re(y_0)$ is a real vector.

Now we shall consider the two situations separately and we first start with the setting of $\| \cdot \|$ being an adjoint complex polytope norm. A fact that $y_0$ is a $2$-strongly unique best approximation for $x$ follows directly from Theorem \ref{twr2strgeneral} applied to a vector $x-y_0 \in \mathbb{R}^n$. To have the setting of adjoint complex polytope norms done, we are left with proving that $y_0$ is not an $\alpha$-strongly unique best approximation for any $\alpha<2$. Indeed, let us assume the contrary and let $f_1, \ldots, f_N: \mathbb{C}^n \to \mathbb{C}$ be the real essential system of facets, i.e.
$$\|z\| = \max \{|f_j(z)|: \ 1 \leq j \leq N\}$$
for $z \in \mathbb{C}^n$. By considering $x-y_0$ instead of $x$ we can suppose that $y_0=0$ and by rescalling we can further assume that $\|x\|=1$. Now let $y \in Y$ be any fixed non-zero vector with real coordinates. Then by the assumed $\alpha$-strong uniqueness for every $t > 0$ we have
$$\|x-tiy\|^{\alpha} \geq 1 + rt^{\alpha},$$
for some constant $r>0$. We note that if $1 \leq j \leq N$ is such that $|f_j(x)|<1$ then for $t>0$ small enough we have $|f_j(x-tiy)|<1$. Hence, for $t>0$ sufficiently small, the norm $\|x-tiy\|$ must be realized for an index $1 \leq j \leq N$ such that $|f_j(x)|=1$. By taking $t=\frac{1}{k}$ for positive integer $k$ and letting $k \to \infty$, we see that some index $1 \leq j \leq N$ will realize the norm $\|x-tiy\|$ infinitely many times. In other words, there exists an index $1 \leq j \leq N$ such that $|f_j(x)|=1$ and for infinitely many $k \geq 1$ we have
$$\left |f_j\left ( x-\frac{i}{k}y \right ) \right |^{\alpha} = \left | 1 - \frac{i}{k}f_j(y) \right |^{\alpha} \geq 1 + \frac{r}{k^{\alpha}}.$$
Obviously we must have $a=f_j(y) \neq 0$. This can be rewritten as
$$k^{\alpha}\left ( \left ( 1 + \frac{a^2}{k^2} \right )^{\frac{\alpha}{2}} - 1) \right ) \geq r.$$
However, by the L'Hospital's rule and the assumption $\alpha<2$ we have
$$\lim_{t \to 0^+} \frac{\left ( 1 + t^2a^2 \right )^{\frac{\alpha}{2}} - 1}{t^{\alpha}} = \lim_{t \to 0^+} \frac{a^2(\left ( 1 + t^2a^2 \right )^{\frac{\alpha}{2}-1}}{t^{\alpha-2}} = \lim_{t \to 0^+} a^2 \frac{t^{2-\alpha}}{\left ( 1 + t^2a^2 \right )^{1-\frac{\alpha}{2}}} =0.$$
Hence, for $t=\frac{1}{k}$ with $k$ sufficiently large we get a contradiction and the conclusion follows.


We can now turn our attention to the case of $\| \cdot \|$ being a complex polytope norm. Let $u_1, \ldots, u_N \in \mathbb{R}^n$ be an essential system of vertices of $\| \cdot \|$, so that for $z \in \mathbb{C}^n$ we have
$$\|z\| = \min \left \{ \sum_{j=1}^{N} |\lambda_j|: \ z = \sum_{j=1}^{N} \lambda_j u_j \right \}.$$
As before we can assume that $y_0=0$ and $\|x_0\|=1$.

We shall prove that for some constant $r>0$ we have
\begin{equation}
\label{ineqteza}
\|x-y\| \geq 1 + r\|y\|^2
\end{equation}
for every $y \in Y$, as squaring this inequality yields the desired $2$-strong uniqueness. Moreover, by Lemma \ref{lemcompact} it is enough to establish \eqref{ineqteza} only for $\|y\| \leq 1$.

Because the norm $ \| \cdot \|$, when considered as a norm on $\mathbb{R}^n$, is a polytope norm (with the unit ball $\conv \{\pm u_1, \ldots, \pm u_N\}$) and $Y$ is a real subspace, we know that $0$ is a $1$-strongly unique best approximation for $x$ in a subspace of real parts of vectors from $Y$. In other words, there exists a constant $r_1>0$ such that for every $y \in Y$ we have
$$\|x-\re(y)\| \geq 1 + r_1\|\re(y)\|.$$
In particular, if $y \in Y$ is such a vector that
$$\|\re(y)\| \geq C \|y\|^2$$
for some constant $C \in (0, 1)$ to be specified later, then \eqref{ineqteza} holds with a constant $Cr_1$. Thus, let us assume that the opposite inequality is true. In particular
$$\|\re(y)\| \leq C \|y|^2 \leq C \|y\|.$$
However, on the other hand we have
$$\|y\| = \|\re(y) + i\im(y)\| \leq \|\re(y)\| + \|\im(y)\|$$
and hence
\begin{equation}
\label{ineqim}
\|\im(y)\| \geq (1-C)\|y\|.
\end{equation}
We can now estimate
$$\|x-y\| = \|x - \re(y) - i\im(y)\| \geq \|x - i \im(y)\| - \|\re(y)\| \geq \|x - i \im(y)\| - C\|y\|^2.$$
If we now prove that $\|x - i \im(y)\| \geq 1 + C'\|y\|^2$ for some constant $C'>C$ then the inequality \eqref{ineqteza} will follow with the constant $C'-C>0$. Let us now take a representation
$$x - i \im(y) = (a_1 - ib_1)u_1 + \ldots + (a_N - ib_N)u_N$$
such that $\|x - i \im(y)\| = \sum_{j=1}^{N} |a_j - ib_j| = \sum_{j=1}^{N} \sqrt{a_j^2 + b_j^2}$. Then as $x, \im(y)$ and $u_1, \ldots, u_N$ are all real vectors we have $x = \sum_{j=1}^{N} a_ju_j$ and $\im(y) = \sum_{j=1}^{N} b_ju_j$ so in particular 
\begin{equation}
\label{oszacowania}
\sum_{j=1}^{n} |a_j| \geq \|x\|=1 \quad \text{  and  } \quad \sum_{j=1}^{n} |b_j| \geq \|\im(y)\|.
\end{equation}
Moreover, we note that since $\|x-i \im(y)\| \leq \|x\| + \|\im(y)\| \leq 1 + \|y\| \leq 2$ (using Lemma \ref{lemreal}) we have $|a_j|, |b_j| \leq 2$ for every $1 \leq j \leq N$. Let us now observe that for any two real numbers $a$, $b$ with $|a|, |b| \leq 2$ the following inequality is true
$$\sqrt{a^2 + b^2} \geq |a| + \frac{1}{8}b^2.$$
Indeed, after squaring this is equivalent to
$$b^2 \geq \frac{1}{4}|a|b^2 + \frac{1}{64}b^4.$$
However, from $|a|, |b| \leq 2$ we get
$$\frac{1}{4}|a|b^2 + \frac{1}{64}b^4 \leq \frac{1}{2}b^2 + \frac{1}{16}b^2 = \frac{9}{16} b^2 \leq b^2.$$
Since $|a_j|, |b_j| \leq 2$ we can use the above estimate, combined with the the Cauchy-Schwarz inequality, and estimates \eqref{ineqim},  \eqref{oszacowania} to obtain
$$\|x - i\im(y)\| = \sum_{j=1}^{N} \sqrt{a_j^2 + b_j^2} \geq \sum_{j=1}^{N} \left ( |a_j| + \frac{1}{8} |b_j|^2 \right ) \geq 1 + \frac{1}{8} \sum_{j=1}^{N} |b_j|^2$$
$$ \geq 1 + \frac{1}{8N} \left ( \sum_{j=1}^{N} |b_j| \right )^2 \geq 1 + \frac{1}{8N} \|\im(y)\|^2 \geq 1 + \frac{1-C}{8N}\|y\|^2.$$
Thus, we can take $C' = \frac{1-C}{8N}$ and it is enough to choose $C>0$ so that $\frac{1-C}{8N} > C$. In particular, if we take $C=\frac{1}{16N+1}$ then $C'=2C$. Altogether, the desired inequality \eqref{ineqteza} holds for $\|y\| \leq 1$ with a constant $ r = \frac{r_1}{16N+1}>0$.
\end{proof}

\end{document}
